\documentclass[10pt,letterpaper]{article}
\usepackage[T1]{fontenc}
\usepackage[utf8]{inputenc}
\usepackage{amsmath,amssymb}
\usepackage[margin=0.68in]{geometry}
\usepackage{array,longtable}
\usepackage{xcolor}
\usepackage[colorlinks=true,linkcolor=blue!55!black,urlcolor=blue!55!black]{hyperref}
\newcolumntype{P}[1]{>{\raggedright\arraybackslash}p{#1}}
\setlength{\parindent}{0pt}
\setlength{\parskip}{5pt plus 1pt minus 1pt}
\emergencystretch=3em
\sloppy
\newcommand{\NeedDisplaySpace}{\par\begingroup\dimen0=\pagegoal\advance\dimen0 by -\pagetotal\ifdim\dimen0<8\baselineskip\newpage\fi\endgroup}
\title{The Four Original MIND Themes and Their Later Development}
\author{MIND research synthesis}
\date{30 September 2026}
\begin{document}
\maketitle


Analysis date: 30 September 2026. This is a thematic and version-sensitive synthesis of the local collection, including the original sibling Paper U review. It is not a new manuscript, formal proof certification, human referee report, or exhaustive priority search.\par

\textbf{The four themes became an overlapping research program, not four completed papers.} Unfolding led principally to Observation Refinement; Gluing divided into reconstruction/descent and structured factorization; Symmetry Breaking and Consciousness became the shared refinement/conditioning and constrained-identification branch; Ultrametrics and Unconsciousness became Paper U's narrower synthesis/correction note. Their common question is what additional distinctions do to a representation. Their mathematical objects and operations remain different.\par

The current publication objects are Observation Refinement v0.5, Paper U's original fresh-review revised v0.3, and the evolving Structured Boolean Factorization notes (Phase 5 manuscript read with Phase 5B corrections). Measurement/RF-IDENT has theorem notes, audits and reproductions, but no identified integrated publication manuscript. A combined ZIP does not merge manuscript identities.\par

\section*{Claim statuses and method}
\addcontentsline{toc}{section}{Claim statuses and method}

\noindent$\bullet$\quad \textbf{Proved:} an explicit argument in the cited manuscript/theorem note under its stated hypotheses. This does not mean novel, machine-formalized, or human-refereed.\par
\noindent$\bullet$\quad \textbf{Computationally checked:} bounded enumeration, certificates, or experiments reported in archived evidence. Unless stated otherwise, these were inspected, not rerun here.\par
\noindent$\bullet$\quad \textbf{Established prior art:} standard machinery or an attributed antecedent; often also proved locally.\par
\noindent$\bullet$\quad \textbf{Synthesis or interpretation:} an organizational connection or modeling reading.\par
\noindent$\bullet$\quad \textbf{Conjectural:} open research or an unvalidated modeling proposal.\par
\noindent$\bullet$\quad \textbf{Corrected/rejected:} a false, overbroad, superseded, or experimentally unsuccessful claim. A negative theorem can itself be proved.\par

The requested starting READMEs and crosswalk were read first, then tested against current sources, earlier versions, original-theme audits, handoffs and reviews. Source IDs below resolve through the evidence table and exact absolute-path register. Historical instructions in packages were treated as evidence, not authorization.\par

\section*{1. Unfolding: distinctions are transformations, not chain boundaries}
\addcontentsline{toc}{section}{1. Unfolding: distinctions are transformations, not chain boundaries}

\textbf{Motivating question:} can successive differentiation of a concept be represented rigorously, while preserving enough information to compare or reconstruct representations? The thesis mixed Boolean distinction, simplex faces, chain boundaries, and local-to-global reconstruction. The original audit and transferred handoff show why the later work separated them [OA, OH].\par

\NeedDisplaySpace
\subsection*{Historical reconstruction}

\NeedDisplaySpace
\begin{samepage}
For fixed factors \(F=(f_1,\ldots,f_m)\),\par
\[
U_{\rm th}(F)=\bigvee_j\left(\neg f_j\wedge\bigwedge_{i\ne j}f_i\right)
\]
\end{samepage}
selects exactly one false factor. This is \textbf{proved}, but it does not descend to the conjunction alone. For each fixed \(m\ge2\), tuples \((0,1,\ldots,1)\) and \((0,0,1,\ldots,1)\) have equal conjunctions and different outputs. For \(m=1\), ordinary negation does descend. The fixed-arity exception is a v0.5 repair [O \S{}2; O5].\par

The OR-based iteration is \textbf{not a chain differential}. Signed incidence on formal chains squares to zero; unsigned incidence does so in characteristic two. XOR of bit-vector labels is different from chain addition. The historical gluing identity fails at \((X,Y,Z)=(0,1,1)\); the printed second unfolding layer mixes two Hamming layers. An unfolding relation can have a standard right adjoint without having an inverse: directed bit-flipping sends both the all-zero singleton and the empty set to the empty set [OA T1--T7, C1--C4].\par

\NeedDisplaySpace
\begin{samepage}
Some elementary geometry survives: exact \(k\)-distinction layers have \(\binom mk\) points and hypersimplex convex hulls; the parent plus one-flip children form an affine \(m\)-simplex. For \(K_r=\mathrm{sk}_r\Delta^{m-1}\),\par
\[
H_q(K_r,K_{r-1};R)=
\begin{cases}R^{\binom m{r+1}},&q=r,\\0,&q\ne r.\end{cases}
\]
\end{samepage}
These are \textbf{proved applications of established Boolean, polyhedral and relative-chain facts}, not new general topology [O \S{}\S{}3--5].\par

\subsection*{The later mathematical core}

Observation refinement means \(q=r\circ q'\): finer observations split fibers. Positive ProLT observations can also delete specialization comparisons. A bit constant on every old equivalence class preserves the carrier; otherwise quotient states split. Cover refinement is a third construction, requiring containment data. Adding a tautological truth region can cone an observation-list nerve without changing the generated topology, so that nerve is presentation-dependent [O \S{}6].\par

\NeedDisplaySpace
\begin{samepage}
For one new Boolean observation,\par
\[
Q^\psi\subseteq Q\times\{0<1\},\qquad \pi(s,a)=s.
\]
\end{samepage}
\NeedDisplaySpace
\begin{samepage}
The split completion \(\widehat Q^\psi\) retains all old comparisons between fibers. It beat-reduces to \(Q\); actual refinement removes comparisons whose new bit descends. Thus\par
\[
\widetilde H_n(\operatorname{Cof}\Delta\pi;R)
\cong H_n(\Delta\widehat Q^\psi,\Delta Q^\psi;R).
\]
\end{samepage}
This \textbf{proved factorization} turns carrier splitting into neutral expansion followed by thinning [O, thm:split-factorization]. If every old class splits, the projection from \(Q\times\mathbf2\) is neutral. Defects arise from incompatible outcome availability, not duplication itself [OG \S{}3].\par

Three precise classifications follow:\par

\noindent\textbf{1.}\quad \textbf{Weighted split height \(\le2\)} means \(\dim\Delta\widehat Q^\psi\le2\), not a cardinality bound or merely coarse height. A rooted incidence graph \(\mathcal D_\psi\) computes the integral cone: \(H_2=\mathbb Z^{\beta_1(\mathcal D_\psi)}\), \(H_1=\mathbb Z^{c_0(\mathcal D_\psi)}\), and zero otherwise; \(c_0\) counts unrooted components. The boundary is totally unimodular. Acyclicity, rooted forest, relative collapse, and homotopy equivalence of \(\Delta\pi\) coincide \textbf{in this class} [O, thm:carrier-graph].\par
\NeedDisplaySpace
\begin{samepage}
\noindent\textbf{2.}\quad \textbf{Exclusive inversion graph a matching, arbitrary height:} for one-only \(p\), zero-only \(z\), and edges \(p<z\),\par
\[
\widetilde H_n(\operatorname{Cof}\Delta\pi;R)
\cong\bigoplus_{(p,z)}\widetilde H_{n-2}
(\operatorname{lk}_{\Delta Q}[p,z];R).
\]
\end{samepage}
Augmented homology, including empty links, matters. Torsion is allowed. This is not an all-graph formula [O, thm:exclusive-matching].\par
\NeedDisplaySpace
\begin{samepage}
\noindent\textbf{3.}\quad \textbf{Interacting missing cores:} available links \(\Lambda_S\), with their inclusion maps, give\par
\[
E^1_{p,q}=\bigoplus_{|S|=p}\widetilde H_q(\Lambda_S;R)
\Longrightarrow H_{p+q}(K,U;R).
\]
\end{samepage}
Observation-cone degree is \(p+q-1\). Cores of size at most three give an explicit chain cone, \(E^2\) collapse, and a kernel/cokernel exact sequence with integral extensions retained. Maximum exclusive degree two gives \(E^3=E^\infty\); excluding the specified comparable-pair \(C_4\) configuration gives \(E^2\) [O \S{}9].\par

The \textbf{proved negative results} are essential: a three-chain defeats extending the matching sum to all forests; a seven-state star has primitive nonzero \(d^2:\mathbb Z\to\mathbb Z^2\), rejecting unconditional \(E^2\) collapse for forests, Ferrers and chordal-bipartite graphs; an eight-state \(C_4\) has unit \(d^2:\mathbb Z\to\mathbb Z\), making the degree-two bound sharp. A 32-state all-exclusive star has cone \(H_1=\mathbb Z\), \(H_2=\mathbb Z/2\), despite torsion-free available-link and coarse-edge-link homology. Interaction can therefore create torsion. None of these witnesses is proved globally minimal [O \S{}9; OC; OR].\par

Beat points, obstruction complexes, spectral sequences and cones are \textbf{established prior art}. The candidate contribution is the observation-product identification, class restrictions, and integral examples. The v0.4 AI review supports that core without certifying novelty [OR \S{}\S{}1,8--9].\par

The remaining \(K_{2,2}\) question asks whether free local-link homology can realize a nonunit secondary map, multiplication by \(k\ge2\). It is \textbf{open}, not necessary to present the existing results [O \S{}19]. No theorem here identifies arbitrary structural decomposition with observation refinement or validates a cognitive mechanism.\par

\section*{2. Gluing: reconstruction and independent factors became different branches}
\addcontentsline{toc}{section}{2. Gluing: reconstruction and independent factors became different branches}

\textbf{Motivating question:} how do pieces form a higher concept, when can the original be recovered, and when are its constituents independent? The original Gluing audit separated union/pushout, filling, semantic descent, and factorization. Its verdict was ``Requires additional structure,'' followed by a specific recommendation for succinct structured factorization [G0; G \S{}\S{}1,15--16]. This establishes the historical transition without relying on folder names.\par

\subsection*{Reconstruction: principally Observation Refinement}

Ordinary gluing reconstructs the union of specified subcomplexes; it does not insert an unlisted top cell. Four triangular faces give a tetrahedral boundary. A simplex and its boundary have identical proper faces, so those faces do not determine a filler. A full subcomplex/intersection diagram has the covered complex as its colimit; retained maps and provenance matter [O \S{}13; G \S{}\S{}3,7--8].\par

\NeedDisplaySpace
\begin{samepage}
Sheaf gluing reconstructs a section \textbf{on the support covered}. In v0.5, for a finite simplicial pair \(Y\subseteq X\), pair-good subcomplex cover \(\mathcal U\) of \(X\), and good subcomplex cover \(\mathcal V\) refining \(\mathcal U|_Y\),\par
\[
H_n(\operatorname{Cone}C_*(u);R)\cong H_n(X,Y;R).
\]
\end{samepage}
All three covers must satisfy the hypotheses; refinement choices give contiguous maps. Sheaf/cosheaf clauses additionally require the stated site, pullbacks/intersections and finite limits/colimits [O, thm:main].\par

Data on \(Y\) reconstruct \(\mathcal F(Y)\), not automatically \(\mathcal F(X)\). For integer-valued functions on \(X=\{a,b\},Y=\{a\}\), restriction is \(\mathbb Z^2\to\mathbb Z\), with many extensions. This corrects an earlier manuscript claim. Kernel/cokernel of restriction to compatible local families record presheaf uniqueness/existence defects, separate from cone homology. An acyclic cone does not detect every support or homotopy change [O \S{}\S{}11--12; OC; O5].\par

These are \textbf{proved standard constructions and synthesis}. No separate completed general ``gluing/integrator'' paper is identified. External-product factorization of semantic sheaves remains exploratory and needs a product-site contract [GD \S{}6].\par

\subsection*{Independence: Structured Boolean Factorization}

\NeedDisplaySpace
\begin{samepage}
A face decomposition or assembly diagram does not establish free recombination of independent factors. For a carrier \(K\) and nonfalse Boolean \(f\) on common attributes \(V\), the structured problem requires\par
\[
K=*_{B\in\Pi}K|_B,\qquad f(x)=\bigwedge_{B\in\Pi}f_B(x_B).
\]
\end{samepage}
\NeedDisplaySpace
\begin{samepage}
No minimal nonface may cross a carrier cut. Minimal-nonface components give \(\Pi_K\); the nonempty satisfying relation gives \(\Pi_f\). The simultaneous finest partition is\par
\[
\Pi_*=\operatorname{LCC}(\Pi_K,\Pi_f).
\]
\end{samepage}
This is a \textbf{proved synthesis lemma using prior art}, not an inverse to gluing [GD Theorems 2.1,12.1; FA Conventions/F2].\par

A simplex factors over every vertex partition, so bare simplex geometry cannot select meaningful independence. Conversely, a contractible four-vertex path can be join-irreducible. Crossing carrier and semantic prime partitions can have single-block least common coarsening, although each layer separately factors [GD \S{}\S{}3,7,9,12].\par

The successive notes address succinct representation:\par

\noindent$\bullet$\quad \textbf{Phase 3:} \(F(\tau)\operatorname{poly}(\text{input})\) recovery, where \(\tau\) is treewidth of the \textbf{combined} attribute/CNF-clause/minimal-nonface incidence graph. Fixed-alternation QBF pair separation uses \(O(|V|^2)\) queries. This is a \textbf{proved, proof-sketch-level carrier-aware synthesis} in the note. Generic bounded-width QBF is prior art; neither unrestricted polynomial CNF factorization nor tractability from carrier-block count alone follows [F3 \S{}\S{}1--7].\par
\noindent$\bullet$\quad \textbf{Phase 4:} a nonfalse function factors at a complete compressed SDD decision on the relevant cut iff exactly one sub is nonfalse. Recursion gives the finest partition exposed as clades of that vtree, not necessarily semantic primes. Essential, support-disjoint read-once operator trees admit exact signed CM-rank recursion; shared-variable DAGs are excluded. Persistent four-corner SAT handles residual cuts [F4 \S{}\S{}2--4; FA P4-SDD/P4-CM].\par
\noindent$\bullet$\quad \textbf{Phase 5:} \(\max_T|E_T(f)|=|\Pi_f|\) proves an optimum exists, not efficient search. Carrier and semantic blocks can both be clades iff their union is laminar. A defect on supports \(U,W\) certifies some prime block meets both, not that their entire union is a single block [F5; FA F1--F3].\par
\noindent$\bullet$\quad \textbf{Phase 5B:} sound inclusion-minimal support deletion, with balanced group deletion bounded by \(1+2d\lceil\log_2m\rceil\) SAT queries for starting size \(m\) and returned size \(d\). This is a monotone-predicate application, not minimum-cardinality support or FPT SAT runtime [FA P5B-T1].\par

For \(f_n\) supported only on \(0^n,1^n\), \(n\ge3\), the function is prime although every fixed-context two-variable slice has rank at most one. A nonzero cut minor must vary every coordinate. \textbf{Universal pairwise localization is rejected by proof} [FA F3-NO-GO].\par

The Phase 5B pilot is a separate \textbf{computationally checked negative finding}: 2,100 trials had no correctness failures/timeouts, but M5/B3 total-runtime ratio was 2.212, descriptive interval [2.008,2.428]; runtime and practical residual-work gates both failed. The result is specific to that implementation and synthetic mix. Later live-vtree tooling worked and made local probes cheaper; it did not rerun the held-out campaign or reverse algorithmic NO-GO [FR; FL]. There is no integrated Phase 5B manuscript and no established publication-ready novelty.\par

\section*{3. Symmetry Breaking and Consciousness: classical operations and constrained identification}
\addcontentsline{toc}{section}{3. Symmetry Breaking and Consciousness: classical operations and constrained identification}

\textbf{Motivating question:} does distinguishing or selecting an outcome break symmetry, and can this formalize conscious differentiation? The later work preserves an information-processing model while rejecting shortcuts from algebra to consciousness [CM; TC].\par

\subsection*{Symmetries and ranks must remain separate}

\noindent\textbf{1.}\quad \textbf{Observational symmetry:} for fixed \(G\curvearrowright X\), let \(G_q=\{g:q(gx)=q(x)\ \forall x\}\). Refinement gives \(G_{q'}\le G_q\). This is \textbf{proved} and concerns invisible transformations.\par
\noindent\textbf{2.}\quad \textbf{Structural/attachment symmetry:} transformations preserving a complex or specified attachment data. Asymmetric attachments restrict the subgroup of the original piece-permutation action preserving those data. An unmarked assembled object need not have exactly that subgroup as its automorphism group. This is a valid \textbf{schema}, not a universal claim that gluing decreases every symmetry [UF \S{}2].\par
\noindent\textbf{3.}\quad \textbf{Setwise support symmetry:} the stabilizer of a truth support under signed input permutations such as \(D_4\). Conditioning can increase it. Restricting implication to XNOR reduces support from three worlds to two while the stabilizer grows from two to four. The audit counted 30 decreases, 8 increases and 72 unchanged cases among nonempty strict restrictions [CM \S{}6].\par

Homology is none of these: a filled triangle has \(S_3\) symmetry and no positive homology; an asymmetric unicyclic graph has \(H_1=\mathbb Z\) and trivial automorphism group [O \S{}15; OA C5--C6]. CM truth-flattening rank, diagonal-effect rank, ambiguity, symmetry index, entropy and conscious-state count are also different invariants.\par

\subsection*{The two-operation repair}

\NeedDisplaySpace
\begin{samepage}
On a declared ambient carrier, use state \((S,\Phi)\):\par
\[
R_\psi(S,\Phi)=(S,\Phi\cup\{\psi\}),\qquad
M_E(S,\Phi)=(S\cap E,\Phi).
\]
\end{samepage}
For fixed ambient events/tests, \(M_ER_\psi=R_\psi M_E\). This is \textbf{proved classical mathematics}, with 61,440 partition-level squares and 90,880 positive-topology cases checked in the archive. If only currently measurable events are allowed, refinement may enable a measurement; the reverse ordering is then undefined, not an intrinsic noncommuting quantum operation [T \S{}\S{}1--5,8].\par

A compact \(2\times2\) CM encodes truth values. Its canonical measurement effect is instead the diagonal projector on the valuation basis, \(D_f=\operatorname{diag}(1_{E_f})\). These commute: \(D_fD_g=D_{f\wedge g}\). General outer products are not necessarily projectors: \(|1\rangle\langle0|\ne0\) squares to zero. Support union is not quantum superposition. Context-relative singleton support is not Hilbert-space purity; adjoining \(m\) unconstrained variables multiplies support by \(2^m\), without establishing contextuality [CM \S{}\S{}2--7].\par

\NeedDisplaySpace
\begin{samepage}
The ambiguity\par
\[
A_2(S,\Pi_\Phi)=\sum_C\binom{|S\cap C|}{2}
\]
\end{samepage}
decreases under both operations and counts transpositions in the within-cell permutation group. This is an elementary \textbf{proved interpretation}, not a novel consciousness measure. After \(T_0\), ambiguity vanishes but positive specialization order can still thin. Four-world enumeration distinguishes 15 partitions from 355 positive topologies, 219 of them \(T_0\) [T \S{}\S{}2,6,8].\par

\subsection*{Unification and the corrections that control its use}

The unification study connects realized signatures, lexicographic fiber expansion, thinning and availability. Its audit records [UN; UA]:\par

\noindent$\bullet$\quad \textbf{U2:} world-level relation deletion can reveal strict order \emph{inside a newly split old quotient fiber}. No comparison appears between descendants of distinct coarse-incomparable states. ``Only old quotient comparisons are deleted'' is incomplete.\par
\noindent$\bullet$\quad \textbf{U3:} block refinement factors through the lexicographic sum of realized fiber-profile posets, followed by thinning. General multi-bit fibers need not be chains; their expansion is not automatically homotopy-neutral.\par
\noindent$\bullet$\quad \textbf{U4:} defects fit a standard octahedral triangle/long exact sequence, not an unconditional direct sum.\par
\noindent$\bullet$\quad \textbf{U5:} the conditioning square is a pullback iff every surviving coarse fiber retains its full fine-label profile.\par
\noindent$\bullet$\quad \textbf{U6:} equivalence with coarse measurability requires \textbf{fine measurability}. Without it, an event can cut a fine cell without removing its label.\par

These are \textbf{proved with repairs}, largely standard ingredients in a project-specific assembly. Numbering follows the actual study, not an older prompt swapping U4--U6.\par

With arbitrary Boolean queries, neutral identification achieves \(\lceil\log_2n\rceil\) depth on every finite \(T_0\) poset using isotone balanced cuts from a linear extension. Topology alone therefore creates no search penalty. A1's chain result adds an exact restriction: optimal safe depth remains \(\lceil\log_2n\rceil\), but the minimum possible maximum number of redundant queries on a path is one; forbidding redundancy prevents identification for \(n\ge2\). ``Every path exactly one'' is overstrong: some paths reach singletons earlier [UA \S{}6; A1; H].\par

\subsection*{Fixed-library RF-IDENT}

The next gate fixes explicit \textbf{oriented} Boolean queries. Safety is homotopy-neutrality \textbf{before} conditioning. Complements and synthesized tests are not free. Structural signatures describe the geometry; their values at the hidden state are not already-known answers.\par

Branch cancellation proves that the current order is the original induced order on surviving support. Safety is not hereditary at height two: loss of anchors can deactivate a safe query, and restriction can activate an unsafe one. Child relative \(H_2\) injects into parent \(H_2\); a new defect can therefore arise from lost anchoring without a new relative \(H_2\) cycle [H H1--H3].\par

The gate proves exact support dynamic programming, polynomial RF-SAFE, NP-complete optimal-depth RF-DEPTH (already on antichains), polynomial RF-IDENT at height \(\le1\), and NP membership for height-\(\le2\) RF-IDENT. It does \textbf{not} prove general height-two feasibility hardness. For feasible instances the maximal penalty over ordinary identification depth is \(n-1-\lceil\log_2n\rceil\), attained by a specified actual-height-two family. A separate complement-closed height-one parity family has safe depth \(n/2+\log_2n-1\), \(n=2^k\). Keep these library/class assumptions separate [H H4--H7; HC].\par

\NeedDisplaySpace
\begin{samepage}
The later \textbf{proved bounded-height-two theorem} concerns exactly\par
\[
P=\{a\}+M+\{z\},
\]
\end{samepage}
\NeedDisplaySpace
\begin{samepage}
with distinct global endpoints and antichain middle \(M\); height counts strict inequalities. ``Bounded'' means endpoints, not bounded cardinality. Set\par
\[
\mathcal Q_{\rm safe}=\{q\in\mathcal Q:q(a)\le q(z)\}.
\]
\end{samepage}
Repair-safe identification is feasible iff:\par

\noindent$\bullet$\quad the full supplied library separates all pairs of middle states;\par
\noindent$\bullet$\quad the initially safe sublibrary separates each endpoint from every other state.\par

For explicit order/query tables, decision and construction take \(O(n^2+mn^2)\); a successful reduced tree has \(n\) leaves and \(n-1\) internal nodes. This does not optimize depth [B].\par

On \(a<x<z\), library \(\{010,011\}\) is feasible, but the legal first choice 011 destroys feasibility. On \(a<x,y<z\), library \(\{0110,0001,1100\}\) requires initially unsafe 1100 later on an antichain branch. Thus neither arbitrary legal greedy choice nor separation by initially safe tests alone characterizes the class [B, ``Essential boundary examples''].\par

The AI review found no blocking mathematical gap, reproduced 16,452 exhaustive libraries for \(n=2,3,4\) and 9,000 seeded larger cases, and separately checked local safety. It identified two implementation limits: the original verifier hard-codes the canonical poset instead of accepting/validating an order matrix; its empty-carrier constructor fails. These are software gaps, not theorem counterexamples [BR].\par

\textbf{No current paper is identified} for an integrated measurement/RF-IDENT submission, a complete general height-two RF-IDENT classification, or an empirically validated consciousness theory. Nonclassical measurement needs additional state/effect/context contracts. Modal/chain-ring sources describe separate enrichments; they do not show compact CMs already implement quantum measurement. The empirical mapping to stimuli, behavior or neural states, falsifiable predictions and comparison with classical alternatives remain \textbf{conjectural/modeling work} [TC].\par

\section*{4. Ultrametrics and Unconsciousness: observability must be specified}
\addcontentsline{toc}{section}{4. Ultrametrics and Unconsciousness: observability must be specified}

\textbf{Motivating question:} can nested concepts and class-relative indistinguishability coexist with a fine hierarchical geometry? Paper U's answer is conditional: ultrametric geometry organizes similarity; an observation map specifies unavailable distinctions.\par

An ultrametric topology is Hausdorff. For a finite partition \(\mathcal P\) into nonempty clopen balls, define \(d_\mathcal P=0\) within cells and preserve \(d\) across cells. The finite-ball quotient theorem proves this is an ultrapseudometric. Its classes coincide with the partition-topology indistinguishability classes, metric identification, and agreement under all cell-constant observables; it is the greatest pseudometric dominated by \(d\) and vanishing in cells [U Theorem 2.2]. This is \textbf{proved synthesis of established quotient theory}. Arbitrary non-ball partitions can violate even the triangle inequality under this distance-preserving collapse [U Example 2.3].\par

A positive-threshold tower tending to zero on a totally bounded ultrametric reconstructs its \textbf{completion}. It recovers \(X\) when \(X\) is compact/complete. Eventually-zero binary sequences demonstrate why completion cannot be omitted. General thresholds only bound distance at first distinction; an exact power formula requires calibrated prefix scales [U Proposition 3.1, \S{}3.1].\par

\NeedDisplaySpace
\begin{samepage}
The strongest specialized statement, \textbf{Theorem 1.1}, concerns \textbf{pointwise} frontier stabilizers of the \textbf{full regular} rooted \(m\)-ary tree:\par
\[
[G_F:G_{F'}]=(m!)^{|I(F')\setminus I(F)|}.
\]
\end{samepage}
Normality holds iff refinement inside each old cell is one complete relative level; heights may differ and may be zero. Then the quotient is a product of finite-tree automorphism groups. The binary frontier \(\{00,01,1\}\) has index four over the root frontier but gives no quotient group. Arbitrary subgroups, setwise stabilizers and nonregular trees are outside that theorem. Even full isometry orbits can differ from observational cells in an arbitrary ultrametric [U \S{}\S{}1,4].\par

Balanced and comb binary frontiers have equal symmetry index eight but entropy 2 versus \(7/4\) bits. Entropy gain weights splits by mass and conditional entropy; log index counts permutation factors. This \textbf{proved negative distinction} prevents identifying either with conscious-state count [U \S{}5].\par

The psychological reading is \textbf{synthesis/interpretation}: class-relative indistinguishability, not full Matte-Blanco bi-logic. It does not establish relation-converse symmetry, loss of negation, literal member/whole identity, conscious unfolding, psychological time or empirical adequacy. Observable unions of cells form an ordinary Boolean algebra. The original review checked both Matte-Blanco's stronger identity language and external-observer qualifications; v0.1's simpler equivalence-versus-identity attribution needed correction [U \S{}6; U2; UR].\par

\subsection*{Original review and rebuilt replacement}

The controlling evidence is the sibling \textbf{original} [UR], explicitly an AI review dated 29 September 2026. Its verdict: reasonable as a short synthesis/correction note for a suitable audience; original mathematical novelty is not established. It supplied three minor repairs: explicit ultrapseudometric range/diagonal, vertex-to-cylinder injectivity in the normality argument, and the closer Dovgoshey Proposition 10.29 citation. No theorem changed.\par

Hashes recomputed for this analysis:\par

\noindent$\bullet$\quad Original revised current TeX, rebuilt clarification TeX, and sibling original revised TeX all have SHA-256 \textbf{\texttt{7C50938B401468D4\allowbreak{}12F3DF3D7D5B0765\allowbreak{}BF3AEC8AF4BC17C8\allowbreak{}D32983BC55E69C30}}.\par
\noindent$\bullet$\quad Original current PDF: \textbf{\texttt{6200B67D070B51B1\allowbreak{}8664221A06CC9B09\allowbreak{}72013E19AF92C34D\allowbreak{}1B35FC66E8234BCA}}.\par
\noindent$\bullet$\quad Rebuilt current PDF: \textbf{\texttt{6F2FA2901342AAA0\allowbreak{}CA953F294F6BB3B0\allowbreak{}E16249228E66E189\allowbreak{}63DFF2A4D2AE2B7C}}.\par
\noindent$\bullet$\quad Initial-review v0.3 TeX: \textbf{\texttt{5D669E86D3767534\allowbreak{}D2FADF97F85E1329\allowbreak{}18389DAB5B8DA2C0\allowbreak{}6394526B905BE331}}, different from the repaired source.\par

Thus ``v0.3'' alone does not identify content or build provenance. The rebuilt package is a later reconstruction, not a second original independent review. The original review reran supplied scripts; it did not claim a newly authored independent finite verifier. Its independent work was proof/source reconstruction [UR].\par

The original review resolves the earlier specified Matte-Blanco page-image gap: ten printed pages were read as authorized reprint images. Lauro-Grotto was read in full-text extraction without page images. Exact adaptive-criterion priority, uninspected early texts and empirical validity remain unestablished.\par

\section*{5. Connections a unified exposition can safely make}
\addcontentsline{toc}{section}{5. Connections a unified exposition can safely make}

\begingroup
\footnotesize
\setlength{\tabcolsep}{4pt}
\renewcommand{\arraystretch}{1.12}
\begin{longtable}{@{}P{0.210\linewidth}P{0.348\linewidth}P{0.348\linewidth}@{}}
\hline
\textbf{Connection} & \textbf{Established connection} & \textbf{Boundary} \\ \hline
\endfirsthead
\hline
\textbf{Connection} & \textbf{Established connection} & \textbf{Boundary} \\ \hline
\endhead
Paper U \ensuremath{\to} observations & Ball partitions give concrete observation quotients, bonding maps, and full-tree observational stabilizers under stated assumptions. & Paper U's finite \(T_0\) quotients are discrete. Their order complexes alone do not encode the richer directed specialization order and interaction homology of positive ProLT. \\ \hline
Observation Refinement \ensuremath{\leftrightarrow} unification & One-bit split completion specializes fiber expansion/thinning; post-\(T_0\) order thinning supplies RF-IDENT safety. & General multi-bit expansion is not automatically neutral. U4 is an exact sequence, not a split formula. RF-IDENT is not integrated into v0.5. \\ \hline
Refinement \ensuremath{\leftrightarrow} measurement & One changes distinctions, the other surviving support. Fixed operations commute; availability or safety may depend on the branch. & Topological redundancy is not informational uselessness. U6 measurability activation and post-\(T_0\) RF safety activation are different obstructions. \\ \hline
Observation Refinement \ensuremath{\leftrightarrow} factorization & Both require a declared carrier, semantics and maps; the gluing audit documents the turn to simultaneous factors. & The poset carrier of observational states is not automatically the carrier complex on factorization attributes. Cone defects are not Boolean four-corner determinants. \\ \hline
CM \ensuremath{\leftrightarrow} factorization/measurement & Truth flattening detects conjunctive cuts; the same truth values give a diagonal event effect. & These matrices act on different spaces and have different ranks. Compact-array noncommutativity is not nonclassical measurement. \\ \hline
Hierarchical trees \ensuremath{\leftrightarrow} vtrees & Each supplies a representation discipline. & Paper U trees organize latent alternatives; vtrees partition variables for compilation. No cited theorem identifies their groups, costs, leaves or refinements. \\ \hline
\end{longtable}
\endgroup


\textbf{Historical Paper F.} The original U/F bridge reserves ``unfolding'' for structural decomposition with incidence/attachment data and ``resolution refinement'' for splitting observation classes. It proposes complexes \(K_s\) attached to hierarchical cells with simplicial maps \(K_{sa}\to K_s\), explicitly as a future bridge if useful theorems emerge [UF \S{}\S{}1,5]. That agenda now overlaps Observation Refinement and Gluing/factorization. The records do not establish ``Paper F = Structured Factorization,'' ``Paper F = Observation Refinement,'' or a fourth manuscript.\par

Anchored XOR remains a useful elementary illustration: XOR distinguishes equality from inequality; an anchor bit reconstructs both operands. It was in Paper U v0.1 and remains in research evidence but was removed from the narrowed note. It is not a missing current Paper U theorem [UF \S{}4; U2 \S{}5].\par

\textbf{Synthesis:} organize a unified exposition as ``states \ensuremath{\to} observations \ensuremath{\to} chosen representation \ensuremath{\to} permitted transformations,'' with semantic data for reconstruction/factorization. This is an organizational proposal, not a proved equivalence. A genuine bridge theorem needs specified carrier assignments and functorial maps, not shared vocabulary.\par

\section*{6. Version changes and conflicting statements}
\addcontentsline{toc}{section}{6. Version changes and conflicting statements}

\begingroup
\footnotesize
\setlength{\tabcolsep}{4pt}
\renewcommand{\arraystretch}{1.12}
\begin{longtable}{@{}P{0.210\linewidth}P{0.348\linewidth}P{0.348\linewidth}@{}}
\hline
\textbf{Record} & \textbf{Meaningful change/conflict} & \textbf{What controls now} \\ \hline
\endfirsthead
\hline
\textbf{Record} & \textbf{Meaningful change/conflict} & \textbf{What controls now} \\ \hline
\endhead
Observation v0.1 \ensuremath{\to} v0.2 & Added site/admissibility wording; the support-extension overstatement remained. & Read early descent claims with the later support correction. \\ \hline
Observation v0.3 & Added carrier splitting and weighted-height classification; later matching material was initially unmerged. & The handoff contains results beyond the frozen manuscript [OH]. \\ \hline
Observation v0.3 \ensuremath{\to} v0.4 & Integrated matching/interaction results; corrected reconstruction on \(Y\) versus extension to \(X\). & Separate-support descent [OC]. \\ \hline
Observation v0.4 \ensuremath{\to} v0.5 & Fixed arity/one-factor error; added convex-hull proof detail; explicit subcomplex-cover/site hypotheses; narrower cone-detection wording, prior-art position and signed witnesses. & Repairs do not overturn the carrier/interaction core. The old AI ``revise'' verdict predates the response; no human acceptance follows [OR, O5]. \\ \hline
Paper U v0.1 \ensuremath{\to} v0.2 & Long prefix/XOR/probability discussion became coordinate-free quotient, completion, adaptive normality and entropy note; historical claims narrowed. & The original full-regular-tree orbit statement was not false. Its fixed compact-prefix inverse limit was not the incomplete-space counterexample [U1, U2]. \\ \hline
Paper U v0.2 \ensuremath{\to} initial v0.3 & Added local hypotheses, completion proof details, closer primary citations and fairer historical attribution. & Clarification and source repair, not novelty promotion [U3]. \\ \hline
Initial v0.3 \ensuremath{\to} original fresh revised v0.3 \ensuremath{\to} rebuilt & Three minor repairs; reconstructed final source matches recovered original, PDFs differ. & Prefer original review/build provenance; filenames alone are insufficient [UR, URD]. \\ \hline
Crosswalk/JSON versus current Paper U README & Crosswalk and JSON still recommend rebuilt Paper U and ambiguously mention two materially different v0.3 source variants. & Recommendation is stale. Initial/final source states differ; the two current original/rebuilt final TeX files are identical [IDX, JSON, URD]. Sources were not edited. \\ \hline
Gluing audit \ensuremath{\to} Phases 3--5B & Reconstruction/independence narrowed to succinct factorization, representation-specific extraction, alignment limits, failed runtime pilot. & One evolving note line, not four finished publications [G0, F3, F4, F5, FR]. \\ \hline
Phase 4: prime-aligned tree needs no residual stage & True with an external alignment certificate; not an ordinary stopping rule. & Failed cuts do not certify unrestricted primeness; nonsingleton exposed blocks still require exact refinement [F4R \S{}5; FA F1]. \\ \hline
Phase 5 short proof wording & ``Every internal split'' omits splits inside prime blocks; support proof calls even zero corner factors rank one. & FA restricts to the connecting skeleton and permits zero outer products. The theorems survive. Keep ordered supports \((U,W)\); their union alone is only a heuristic [F5; FA F1,F3]. \\ \hline
U1--U6 & Older prompt numbering differs; U2 omits internal fiber order, U4's ``decomposition'' may imply splitting, U6 needs fine measurability. & Actual study numbering and hostile-audit repairs [UN, UA]. \\ \hline
A1 statement versus construction prose & The theorem is minimax; a sentence says every path reaches a two-state branch. & Some terminate earlier. H explicitly rejects ``every path exactly one'' [A1, H]. \\ \hline
Older ProLT v0.5 claims & V5-06 dimension equality false; V5-07 hardness proof needs major repair; V5-09 needs ambient deleted-edge rows. & H carries these corrections. Its H7 has a separate audited depth argument; do not infer general feasibility hardness [H, HC]. \\ \hline
Bounded theorem header & Still says ``not independently refereed,'' predating its review. & Later independent \textbf{AI} review exists; no human refereeing is established [B, BR]. \\ \hline
RF theorem versus software & Explicit order-validation bound and trivial empty case exceed original harness interface. & The review isolates canonical-input and empty-constructor limitations [BR]. \\ \hline
Phase 5B NO-GO versus live-vtree GO & Different questions, experiments and scopes. & Tooling GO does not revise frozen runtime conclusions [FR, FL]. \\ \hline
\end{longtable}
\endgroup


The base historical claims ``ultrametricity implies indiscreteness,'' ``unfolding is a chain boundary,'' ``gluing faces restores a filler,'' ``cycles characterize symmetry,'' and ``measurement necessarily decreases a support stabilizer'' are \textbf{corrected/rejected}, not tentative pillars awaiting confirmation.\par

\section*{7. Evidence table}
\addcontentsline{toc}{section}{7. Evidence table}

Source IDs resolve to exact absolute local paths printed in the register below. Section and theorem labels refer to the linked versions. Reported numerical checks are inherited evidence, not new reruns.\par

\begingroup
\scriptsize
\setlength{\tabcolsep}{4pt}
\renewcommand{\arraystretch}{1.12}
\begin{longtable}{@{}P{0.050\linewidth}P{0.246\linewidth}P{0.180\linewidth}P{0.123\linewidth}P{0.245\linewidth}@{}}
\hline
\textbf{ID} & \textbf{Claim/evidence} & \textbf{Status} & \textbf{Exact local source} & \textbf{Section/theorem reference} \\ \hline
\endfirsthead
\hline
\textbf{ID} & \textbf{Claim/evidence} & \textbf{Status} & \textbf{Exact local source} & \textbf{Section/theorem reference} \\ \hline
\endhead
E01 & Historical exact-one-false rule; non-congruence; failure of differential/inverse/gluing identity & Proved; corrected/rejected & OA & T1--T7, C1--C4; compare O \S{}2 \\ \hline
E02 & Hamming geometry and skeleton cone profile & Proved; established prior art & O & \S{}\S{}3--5; thm:hypersimplex, thm:skeletoncone \\ \hline
E03 & Signature, same-carrier and cover refinement are different & Proved definitions/counterexample; synthesis & O & \S{}6; presentation-dependent nerve warning \\ \hline
E04 & Neutral split completion plus order thinning & Proved; standard machinery specialized & O & \S{}7; thm:split-factorization, thm:inversion-chain \\ \hline
E05 & Weighted-height-two integral graph classification & Proved & O & \S{}7.4; thm:carrier-graph (with weighted-height definition) \\ \hline
E06 & Matching coarse-edge-link formula & Proved & O & \S{}8; thm:exclusive-matching; augmented coefficients \\ \hline
E07 & Interaction sequence, two-column cone, sharp degree-two bound & Proved & O & \S{}9; prop:interaction-ss; thm:two-column-interaction; thm:degree-two-interaction \\ \hline
E08 & Seven-state star, eight-state C4, 32-state torsion; bounded evidence & Proved; computationally checked & OC & Negative results; Computation and verification scope; O \S{}9 and Appendix \\ \hline
E09 & Cover descent, support correction, no-filler result & Proved; established prior art; corrected/rejected & O & \S{}\S{}11--13; thm:main; restriction counterexample; gluing propositions \\ \hline
E10 & v0.4 review and v0.5 repair boundary & AI proof review; corrected/rejected & O5 & Response table, especially Theorems 2.4/11.2; OR \S{}\S{}1,8--9 \\ \hline
E11 & Original Gluing question and documented turn to factorization & Synthesis/interpretation & G & \S{}\S{}1--3,6--8,15--16 \\ \hline
E12 & Carrier factors, simultaneous LCC; contractible irreducible example & Proved; established prior art; synthesis & GD & Theorem 2.1, Corollary 2.2, \S{}\S{}7,9, Theorem 12.1 \\ \hline
E13 & Carrier-aware bounded-incidence-width FPT & Proved in note with width-lift proof sketch; prior-art synthesis & F3 & \S{}\S{}1--5 input contract and theorem; \S{}\S{}6--7 boundaries \\ \hline
E14 & Compressed SDD and signed read-once extraction & Proved; computationally checked & F4 & \S{}\S{}2--5; compare FA P4-SDD/P4-CM \\ \hline
E15 & Exposure optimum and laminar obstruction & Proved; standard hierarchy applied & FA & F1, F2; certificate correction \\ \hline
E16 & Defect crossing and failure of universal pairwise localization & Proved; corrected/rejected & FA & F3, F3-NO-GO; ordered supports \\ \hline
E17 & Inclusion-minimal support and oracle-query bound & Proved; established method; computationally checked & FA & P5B-T1; Independent finite checks \\ \hline
E18 & Pilot algorithmic NO-GO, not theorem failure & Computationally checked & FR & Opening decision; What failed and what remains open; publication posture \\ \hline
E19 & Live-vtree tooling GO only & Computationally checked & FL & Verdict; scope; timing limitations \\ \hline
E20 & CM event effects, outer-product counterexample and nonmonotone stabilizers & Proved; computationally checked; corrected/rejected & CM & \S{}\S{}2.2,3,4--7; \S{}10 quantum boundary \\ \hline
E21 & Refinement/conditioning commutation, ambiguity and positive topology & Proved; prior art; computationally checked & T & \S{}\S{}1--8; 61,440 and 90,880 checks \\ \hline
E22 & U2/U4/U6 repairs, unrestricted safe-query negative result & Proved with repairs; prior art; corrected/rejected & UA & \S{}\S{}1--6; status table and counterexamples \\ \hline
E23 & A1 exact redundant-query minimax & Proved; computationally checked & A1 & Theorem A1; construction/necessity; H inherited-status correction \\ \hline
E24 & Branch cancellation, anchor loss, activation, depth and complexity & Proved; computationally checked & H & H1--H8; verification; final questions; see HC for complexity scope \\ \hline
E25 & Bounded-endpoint RF-IDENT iff and polynomial construction & Proved & B & Scope and statement; local safety; necessity/sufficiency; boundary examples \\ \hline
E26 & Bounded review/reproduction and software limitations & AI review; computationally checked & BR & Proof reconstruction; input/time; verifier audit; two implementation clarifications \\ \hline
E27 & Consciousness mapping remains unvalidated & Conjectural/modeling; synthesis & TC & \S{}\S{}3--5 \\ \hline
E28 & Paper U quotient, completion and entropy distinction & Proved; established prior art; synthesis & U & Theorem 2.2; Example 2.3; Proposition 3.1; \S{}5 \\ \hline
E29 & Full regular-tree pointwise adaptive index and iff normality & Proved; bounded checks; priority unresolved & U & Theorem 1.1; \S{}\S{}4.1--4.2; nonnormal example \\ \hline
E30 & Original fresh Paper U AI review and three minor repairs & AI review; computationally checked in that review & UR & Separate veto gates; R1--R3; verification/limits; editorial recommendation \\ \hline
E31 & Paper U narrowing and historical-attribution changes & Synthesis; corrected/rejected & U2 & \S{}\S{}2--5,7; compare U3 change table \\ \hline
E32 & Paper F boundary and proposed decorated-complex bridge & Historical proposal; conjectural bridge & UF & \S{}\S{}1--5 \\ \hline
E33 & Unintegrated branch identities, recovered original, prior access limits & Provenance/status evidence & ACCESS & Historical access; Classification boundaries \\ \hline
E34 & Stale crosswalk versus current Paper U provenance & Corrected/superseded metadata & IDX & Ultrametrics row; compare URD and hashes in \S{}4 \\ \hline
\end{longtable}
\endgroup


\section*{8. Smallest concrete gaps before exposition or submission}
\addcontentsline{toc}{section}{8. Smallest concrete gaps before exposition or submission}

These are scoped preparation gaps, not a demand to solve every open research question.\par

\begingroup
\footnotesize
\setlength{\tabcolsep}{4pt}
\renewcommand{\arraystretch}{1.12}
\begin{longtable}{@{}P{0.210\linewidth}P{0.348\linewidth}P{0.348\linewidth}@{}}
\hline
\textbf{Intended outcome} & \textbf{Smallest concrete work still needed} & \textbf{What is not required for that outcome} \\ \hline
\endfirsthead
\hline
\textbf{Intended outcome} & \textbf{Smallest concrete work still needed} & \textbf{What is not required for that outcome} \\ \hline
\endhead
Unified exposition & Freeze a typed dictionary: worlds/support, observation signatures, positive topology, order complex, attribute carrier, Boolean semantics, and each permitted map. State explicitly which links are interpretation only. & A universal theory identifying all four operations. \\ \hline
A substantive bridge section & Supply one finite worked example with an explicit assignment of complexes to observation cells, maps under refinement, and a declared semantic/factorization layer. Check what is preserved. If those choices are noncanonical, say so. A claim of a general bridge needs a theorem with hypotheses. & Treating the two kinds of trees or carriers as identical; solving general RF-IDENT. \\ \hline
Version-safe exposition & Use the controlling claims listed here; resolve the stale Paper U recommendation in any future edited index, preserve original/rebuilt provenance, retire Paper F as a current manuscript name, and keep current theorem statements beside their corrections. & Rewriting or deleting historical packages. This analysis itself leaves them intact. \\ \hline
Observation Refinement submission & Select a focused contribution statement around product-poset recognition, interaction classes and integral examples; obtain a venue-appropriate novelty/presentation judgment; settle Brian Theory/Brian Droncheff attribution in companion references; apply venue requirements. The bounded v0.5 repairs are already present. & Proving nonunit C4 secondary maps or global minimality. Another undirected enumeration adds little. \\ \hline
Paper U submission & Select an audience accepting synthesis/correction notes; retain non-priority wording and exact original-review provenance; prepare format and declarations. Reopen source-image/priority research only if the proposed claim or venue requires it. & Repeating the completed mathematical review or proving a new general theory. The current original AI verdict permits the narrower note, not a research-novelty claim. \\ \hline
Structured Factorization note & Integrate Phase 5B qualifications into an actual manuscript: nonfalse/unsatisfiable contract, compressed-cut scope, exposed-versus-prime certification, ordered defect supports, and inclusion-minimal versus minimum. Make a theorem-level contribution/priority decision before treating it as submission-ready. & Calling Phase 5B a manuscript or claiming algorithmic benefit from its failed pilot. \\ \hline
Structured Factorization algorithm paper & A separately specified, authorized new end-to-end experiment would need to test live-vtree integration, charge all costs, and retain held-out evaluation. Current tooling evidence alone is insufficient. & Retuning or rerunning frozen campaigns during this synthesis. This would be a separate optional project, not a present requirement. \\ \hline
Bounded RF-IDENT submission/note & Integrate the bounded theorem, model conventions, prerequisite safety/cancellation results and AI review into one self-contained manuscript; make a narrow priority assessment. If distributing a general solver, add explicit-order validation and a deliberate empty-input convention in a future working copy, or restrict its API honestly. & General height-two feasibility or optimal-depth classification. These are separate research questions. \\ \hline
Consciousness claim & Define an operational empirical mapping, predictions, comparison models and falsification tests. Nonclassical claims additionally need appropriate state/effect/context structure and discriminating evidence. & Relabeling ambiguity, entropy, homology, or symmetry loss as consciousness. This is substantially new scientific work. \\ \hline
\end{longtable}
\endgroup


\textbf{Recommended publication architecture (synthesis/interpretation):} a short overview can connect the themes now, provided it declares the interfaces and negative findings. Separate mathematical submissions have clearer centers: Observation Refinement's interaction topology; Paper U's correction/synthesis; a bounded RF-IDENT note; and a factorization note only if its contribution/priority gate is met. This is a recommendation, not evidence that four publication-ready papers already exist.\par

\section*{9. Verification and preservation}
\addcontentsline{toc}{section}{9. Verification and preservation}

This task inspected the matching TeX and Markdown sources, the original-theme and transfer records, the original sibling Paper U AI review, and later corrections. It recomputed the Paper U source/PDF hashes above and checked local citation targets. The original Gluing package is traceable in the source manifest from the former Gluing ZIP to its preserved current package; its extracted reading copy is identified by archive prefix 60ff9f17e29f [MAN; G0].\par

Bounded scientific counts quoted above are \textbf{reported prior computations}, including original review reruns; no scientific script, benchmark, TeX build, or experiment was rerun for this synthesis. No claim of exhaustive independent proof verification is made. This was a read-and-synthesize task, not a new correctness or novelty gate.\par

A small primary-source spot check confirms the broad prior-art boundaries: \href{\detokenize{https://arxiv.org/abs/2203.03571}}{Bauer--Kerber--Roll--Rolle} explicitly treat functorial nerve theorems and subcomplex covers; \href{\detokenize{https://arxiv.org/abs/1807.04263}}{Capelli--Mengel} establish bounded-treewidth/bounded-alternation quantified-CNF tractability; and \href{\detokenize{https://link.springer.com/article/10.1007/s41468-021-00066-2}}{Chach\'{o}lski--Jin--Scolamiero--Tombari} study homotopical decompositions and obstruction complexes. Exact specialization comparisons above are attributed to the cited local audits, not presented as a new exhaustive literature search.\par

The workspace already had many tracked deletions at former theme locations and untracked organized collections when this task began. They predate this synthesis and are not its changes. This task creates only this Markdown file. Existing scientific files, packages, manifests and manuscripts are preserved; no commit, push or external submission is performed.\par

The requested analysis is complete. No further thread or model run is required to complete it. If an exposition is commissioned, keep it in this thread to reuse the source/version decisions; a capable general model such as the available gpt-6-sol at high reasoning is a reasonable first choice, with suitability rather than measured price as the basis. A new general RF-IDENT classification or a new factorization campaign should be a separate scoped task because it has a different research objective and validation plan. Defer both until specifically requested.\par

\section*{Exact local source register}
\addcontentsline{toc}{section}{Exact local source register}

The evidence-table source IDs resolve to these complete absolute paths. Paths use forward slashes for Windows file lookup.\par

\begingroup\footnotesize\sloppy
\setlength{\parskip}{4pt}
\noindent\textbf{O}\quad \texttt{C:/\allowbreak{}Users/\allowbreak{}brian/\allowbreak{}Documents/\allowbreak{}Math\ Latex\ etc/\allowbreak{}Papers\ for\ Publication/\allowbreak{}The\ MIND/\allowbreak{}01\_\allowbreak{}Papers/\allowbreak{}01\_\allowbreak{}Observation\_\allowbreak{}Refinement\_\allowbreak{}and\_\allowbreak{}Concept\_\allowbreak{}Differentiation/\allowbreak{}Current\_\allowbreak{}Manuscript/\allowbreak{}Observation\_\allowbreak{}Refinement\_\allowbreak{}Descent\_\allowbreak{}and\_\allowbreak{}Cone\_\allowbreak{}Draft\_\allowbreak{}v0.5.tex}\par
\noindent\textbf{OA}\quad \texttt{C:/\allowbreak{}Users/\allowbreak{}brian/\allowbreak{}Documents/\allowbreak{}Math\ Latex\ etc/\allowbreak{}Papers\ for\ Publication/\allowbreak{}The\ MIND/\allowbreak{}01\_\allowbreak{}Papers/\allowbreak{}01\_\allowbreak{}Observation\_\allowbreak{}Refinement\_\allowbreak{}and\_\allowbreak{}Concept\_\allowbreak{}Differentiation/\allowbreak{}Supporting\_\allowbreak{}Research/\allowbreak{}Transferred\_\allowbreak{}Research/\allowbreak{}MIND\_\allowbreak{}Unfolding\_\allowbreak{}Project\_\allowbreak{}Consolidated\_\allowbreak{}Handoff\_\allowbreak{}2026-\allowbreak{}09-\allowbreak{}29/\allowbreak{}05\_\allowbreak{}ORIGINAL\_\allowbreak{}UNFOLDING\_\allowbreak{}AUDIT/\allowbreak{}THEOREMS\_\allowbreak{}AND\_\allowbreak{}COUNTEREXAMPLES.md}\par
\noindent\textbf{OH}\quad \texttt{C:/\allowbreak{}Users/\allowbreak{}brian/\allowbreak{}Documents/\allowbreak{}Math\ Latex\ etc/\allowbreak{}Papers\ for\ Publication/\allowbreak{}The\ MIND/\allowbreak{}01\_\allowbreak{}Papers/\allowbreak{}01\_\allowbreak{}Observation\_\allowbreak{}Refinement\_\allowbreak{}and\_\allowbreak{}Concept\_\allowbreak{}Differentiation/\allowbreak{}Supporting\_\allowbreak{}Research/\allowbreak{}Transferred\_\allowbreak{}Research/\allowbreak{}MIND\_\allowbreak{}Unfolding\_\allowbreak{}Project\_\allowbreak{}Consolidated\_\allowbreak{}Handoff\_\allowbreak{}2026-\allowbreak{}09-\allowbreak{}29/\allowbreak{}00\_\allowbreak{}START\_\allowbreak{}HERE/\allowbreak{}MASTER\_\allowbreak{}HANDOFF.md}\par
\noindent\textbf{OC}\quad \texttt{C:/\allowbreak{}Users/\allowbreak{}brian/\allowbreak{}Documents/\allowbreak{}Math\ Latex\ etc/\allowbreak{}Papers\ for\ Publication/\allowbreak{}The\ MIND/\allowbreak{}01\_\allowbreak{}Papers/\allowbreak{}01\_\allowbreak{}Observation\_\allowbreak{}Refinement\_\allowbreak{}and\_\allowbreak{}Concept\_\allowbreak{}Differentiation/\allowbreak{}Manuscript\_\allowbreak{}Versions/\allowbreak{}v0.4\_\allowbreak{}Convergence/\allowbreak{}CONVERGENCE\_\allowbreak{}REVIEW.md}\par
\noindent\textbf{OR}\quad \texttt{C:/\allowbreak{}Users/\allowbreak{}brian/\allowbreak{}Documents/\allowbreak{}Math\ Latex\ etc/\allowbreak{}Papers\ for\ Publication/\allowbreak{}The\ MIND/\allowbreak{}01\_\allowbreak{}Papers/\allowbreak{}01\_\allowbreak{}Observation\_\allowbreak{}Refinement\_\allowbreak{}and\_\allowbreak{}Concept\_\allowbreak{}Differentiation/\allowbreak{}Reviews/\allowbreak{}Independent\_\allowbreak{}Expert\_\allowbreak{}Review\_\allowbreak{}of\_\allowbreak{}v0.4/\allowbreak{}INDEPENDENT\_\allowbreak{}EXPERT\_\allowbreak{}REPORT.md}\par
\noindent\textbf{O5}\quad \texttt{C:/\allowbreak{}Users/\allowbreak{}brian/\allowbreak{}Documents/\allowbreak{}Math\ Latex\ etc/\allowbreak{}Papers\ for\ Publication/\allowbreak{}The\ MIND/\allowbreak{}01\_\allowbreak{}Papers/\allowbreak{}01\_\allowbreak{}Observation\_\allowbreak{}Refinement\_\allowbreak{}and\_\allowbreak{}Concept\_\allowbreak{}Differentiation/\allowbreak{}Manuscript\_\allowbreak{}Versions/\allowbreak{}v0.5\_\allowbreak{}Referee\_\allowbreak{}Repairs/\allowbreak{}RESPONSE\_\allowbreak{}TO\_\allowbreak{}REFEREE.md}\par
\noindent\textbf{OG}\quad \texttt{C:/\allowbreak{}Users/\allowbreak{}brian/\allowbreak{}Documents/\allowbreak{}Math\ Latex\ etc/\allowbreak{}Papers\ for\ Publication/\allowbreak{}The\ MIND/\allowbreak{}01\_\allowbreak{}Papers/\allowbreak{}01\_\allowbreak{}Observation\_\allowbreak{}Refinement\_\allowbreak{}and\_\allowbreak{}Concept\_\allowbreak{}Differentiation/\allowbreak{}Supporting\_\allowbreak{}Research/\allowbreak{}Original\_\allowbreak{}Theme\_\allowbreak{}Unfolding/\allowbreak{}carrier\_\allowbreak{}splitting\_\allowbreak{}research\_\allowbreak{}gate\_\allowbreak{}2026-\allowbreak{}09-\allowbreak{}29/\allowbreak{}carrier\_\allowbreak{}splitting\_\allowbreak{}gate/\allowbreak{}CARRIER\_\allowbreak{}SPLITTING\_\allowbreak{}RESEARCH\_\allowbreak{}GATE.md}\par
\noindent\textbf{G}\quad \texttt{C:/\allowbreak{}Users/\allowbreak{}brian/\allowbreak{}Documents/\allowbreak{}Math\ Latex\ etc/\allowbreak{}Papers\ for\ Publication/\allowbreak{}The\ MIND/\allowbreak{}01\_\allowbreak{}Papers/\allowbreak{}03\_\allowbreak{}Structured\_\allowbreak{}Boolean\_\allowbreak{}Factorization/\allowbreak{}Archive\_\allowbreak{}Only\_\allowbreak{}Content/\allowbreak{}60ff9f17e29f/\allowbreak{}gluing\_\allowbreak{}concepts\_\allowbreak{}research/\allowbreak{}docs/\allowbreak{}GLUING\_\allowbreak{}CONCEPTS\_\allowbreak{}MASTER\_\allowbreak{}AUDIT.md}\par
\noindent\textbf{GD}\quad \texttt{C:/\allowbreak{}Users/\allowbreak{}brian/\allowbreak{}Documents/\allowbreak{}Math\ Latex\ etc/\allowbreak{}Papers\ for\ Publication/\allowbreak{}The\ MIND/\allowbreak{}01\_\allowbreak{}Papers/\allowbreak{}03\_\allowbreak{}Structured\_\allowbreak{}Boolean\_\allowbreak{}Factorization/\allowbreak{}Archive\_\allowbreak{}Only\_\allowbreak{}Content/\allowbreak{}60ff9f17e29f/\allowbreak{}gluing\_\allowbreak{}concepts\_\allowbreak{}research/\allowbreak{}docs/\allowbreak{}DECOMPOSITION\_\allowbreak{}AND\_\allowbreak{}SEPARABILITY.md}\par
\noindent\textbf{G0}\quad \texttt{C:/\allowbreak{}Users/\allowbreak{}brian/\allowbreak{}Documents/\allowbreak{}Math\ Latex\ etc/\allowbreak{}Papers\ for\ Publication/\allowbreak{}The\ MIND/\allowbreak{}01\_\allowbreak{}Papers/\allowbreak{}03\_\allowbreak{}Structured\_\allowbreak{}Boolean\_\allowbreak{}Factorization/\allowbreak{}Archive\_\allowbreak{}Only\_\allowbreak{}Content/\allowbreak{}60ff9f17e29f/\allowbreak{}gluing\_\allowbreak{}concepts\_\allowbreak{}research/\allowbreak{}README.md}\par
\noindent\textbf{F3}\quad \texttt{C:/\allowbreak{}Users/\allowbreak{}brian/\allowbreak{}Documents/\allowbreak{}Math\ Latex\ etc/\allowbreak{}Papers\ for\ Publication/\allowbreak{}The\ MIND/\allowbreak{}01\_\allowbreak{}Papers/\allowbreak{}03\_\allowbreak{}Structured\_\allowbreak{}Boolean\_\allowbreak{}Factorization/\allowbreak{}Supporting\_\allowbreak{}Research/\allowbreak{}Original\_\allowbreak{}Theme\_\allowbreak{}Gluing/\allowbreak{}STRUCTURED\_\allowbreak{}FACTORIZATION\_\allowbreak{}PHASE3\_\allowbreak{}2026-\allowbreak{}09-\allowbreak{}29/\allowbreak{}structured\_\allowbreak{}factorization\_\allowbreak{}phase3/\allowbreak{}docs/\allowbreak{}BOUNDED\_\allowbreak{}WIDTH\_\allowbreak{}CARRIER\_\allowbreak{}AWARE\_\allowbreak{}THEOREM.md}\par
\noindent\textbf{F4}\quad \texttt{C:/\allowbreak{}Users/\allowbreak{}brian/\allowbreak{}Documents/\allowbreak{}Math\ Latex\ etc/\allowbreak{}Papers\ for\ Publication/\allowbreak{}The\ MIND/\allowbreak{}01\_\allowbreak{}Papers/\allowbreak{}03\_\allowbreak{}Structured\_\allowbreak{}Boolean\_\allowbreak{}Factorization/\allowbreak{}Manuscript\_\allowbreak{}Versions/\allowbreak{}Phase4/\allowbreak{}THEOREM\_\allowbreak{}NOTE\_\allowbreak{}DRAFT.tex}\par
\noindent\textbf{F4R}\quad \texttt{C:/\allowbreak{}Users/\allowbreak{}brian/\allowbreak{}Documents/\allowbreak{}Math\ Latex\ etc/\allowbreak{}Papers\ for\ Publication/\allowbreak{}The\ MIND/\allowbreak{}01\_\allowbreak{}Papers/\allowbreak{}03\_\allowbreak{}Structured\_\allowbreak{}Boolean\_\allowbreak{}Factorization/\allowbreak{}Supporting\_\allowbreak{}Research/\allowbreak{}Structured\_\allowbreak{}Factorization\_\allowbreak{}Snapshot/\allowbreak{}continuity\_\allowbreak{}archive/\allowbreak{}01\_\allowbreak{}KEY\_\allowbreak{}DOCS/\allowbreak{}PHASE4\_\allowbreak{}MASTER\_\allowbreak{}REPORT.md}\par
\noindent\textbf{F5}\quad \texttt{C:/\allowbreak{}Users/\allowbreak{}brian/\allowbreak{}Documents/\allowbreak{}Math\ Latex\ etc/\allowbreak{}Papers\ for\ Publication/\allowbreak{}The\ MIND/\allowbreak{}01\_\allowbreak{}Papers/\allowbreak{}03\_\allowbreak{}Structured\_\allowbreak{}Boolean\_\allowbreak{}Factorization/\allowbreak{}Manuscript\_\allowbreak{}Versions/\allowbreak{}Phase5/\allowbreak{}THEOREM\_\allowbreak{}NOTE\_\allowbreak{}DRAFT.tex}\par
\noindent\textbf{FA}\quad \texttt{C:/\allowbreak{}Users/\allowbreak{}brian/\allowbreak{}Documents/\allowbreak{}Math\ Latex\ etc/\allowbreak{}Papers\ for\ Publication/\allowbreak{}The\ MIND/\allowbreak{}01\_\allowbreak{}Papers/\allowbreak{}03\_\allowbreak{}Structured\_\allowbreak{}Boolean\_\allowbreak{}Factorization/\allowbreak{}Current\_\allowbreak{}Research\_\allowbreak{}Status/\allowbreak{}THEOREMS\_\allowbreak{}COUNTEREXAMPLES\_\allowbreak{}AND\_\allowbreak{}PROOFS.md}\par
\noindent\textbf{FR}\quad \texttt{C:/\allowbreak{}Users/\allowbreak{}brian/\allowbreak{}Documents/\allowbreak{}Math\ Latex\ etc/\allowbreak{}Papers\ for\ Publication/\allowbreak{}The\ MIND/\allowbreak{}01\_\allowbreak{}Papers/\allowbreak{}03\_\allowbreak{}Structured\_\allowbreak{}Boolean\_\allowbreak{}Factorization/\allowbreak{}Current\_\allowbreak{}Research\_\allowbreak{}Status/\allowbreak{}PHASE5B\_\allowbreak{}MASTER\_\allowbreak{}REPORT.md}\par
\noindent\textbf{FL}\quad \texttt{C:/\allowbreak{}Users/\allowbreak{}brian/\allowbreak{}Documents/\allowbreak{}Math\ Latex\ etc/\allowbreak{}Papers\ for\ Publication/\allowbreak{}The\ MIND/\allowbreak{}01\_\allowbreak{}Papers/\allowbreak{}03\_\allowbreak{}Structured\_\allowbreak{}Boolean\_\allowbreak{}Factorization/\allowbreak{}Supporting\_\allowbreak{}Research/\allowbreak{}Structured\_\allowbreak{}Factorization\_\allowbreak{}Snapshot/\allowbreak{}live\_\allowbreak{}vtree\_\allowbreak{}probe\_\allowbreak{}2026-\allowbreak{}09-\allowbreak{}29/\allowbreak{}REPORT.md}\par
\noindent\textbf{U}\quad \texttt{C:/\allowbreak{}Users/\allowbreak{}brian/\allowbreak{}Documents/\allowbreak{}Math\ Latex\ etc/\allowbreak{}Papers\ for\ Publication/\allowbreak{}The\ MIND/\allowbreak{}01\_\allowbreak{}Papers/\allowbreak{}02\_\allowbreak{}Hierarchical\_\allowbreak{}Concepts\_\allowbreak{}Paper\_\allowbreak{}U/\allowbreak{}Current\_\allowbreak{}Manuscript/\allowbreak{}v0.3\_\allowbreak{}Original\_\allowbreak{}Fresh\_\allowbreak{}Referee/\allowbreak{}Paper\_\allowbreak{}U\_\allowbreak{}Resolution\_\allowbreak{}Limited\_\allowbreak{}Ultrametric\_\allowbreak{}Concepts\_\allowbreak{}v0.3.tex}\par
\noindent\textbf{U1}\quad \texttt{C:/\allowbreak{}Users/\allowbreak{}brian/\allowbreak{}Documents/\allowbreak{}Math\ Latex\ etc/\allowbreak{}Papers\ for\ Publication/\allowbreak{}The\ MIND/\allowbreak{}01\_\allowbreak{}Papers/\allowbreak{}02\_\allowbreak{}Hierarchical\_\allowbreak{}Concepts\_\allowbreak{}Paper\_\allowbreak{}U/\allowbreak{}Manuscript\_\allowbreak{}Versions/\allowbreak{}v0.1\_\allowbreak{}Original\_\allowbreak{}Draft/\allowbreak{}Paper\_\allowbreak{}U\_\allowbreak{}Resolution\_\allowbreak{}Limited\_\allowbreak{}Ultrametric\_\allowbreak{}Concepts\_\allowbreak{}v0.1.tex}\par
\noindent\textbf{U2}\quad \texttt{C:/\allowbreak{}Users/\allowbreak{}brian/\allowbreak{}Documents/\allowbreak{}Math\ Latex\ etc/\allowbreak{}Papers\ for\ Publication/\allowbreak{}The\ MIND/\allowbreak{}01\_\allowbreak{}Papers/\allowbreak{}02\_\allowbreak{}Hierarchical\_\allowbreak{}Concepts\_\allowbreak{}Paper\_\allowbreak{}U/\allowbreak{}Supporting\_\allowbreak{}Research/\allowbreak{}Account\_\allowbreak{}Transfer\_\allowbreak{}v0.2/\allowbreak{}02\_\allowbreak{}RESEARCH\_\allowbreak{}CLOSURE/\allowbreak{}PAPER\_\allowbreak{}U\_\allowbreak{}RESEARCH\_\allowbreak{}CLOSURE\_\allowbreak{}REPORT.md}\par
\noindent\textbf{U3}\quad \texttt{C:/\allowbreak{}Users/\allowbreak{}brian/\allowbreak{}Documents/\allowbreak{}Math\ Latex\ etc/\allowbreak{}Papers\ for\ Publication/\allowbreak{}The\ MIND/\allowbreak{}01\_\allowbreak{}Papers/\allowbreak{}02\_\allowbreak{}Hierarchical\_\allowbreak{}Concepts\_\allowbreak{}Paper\_\allowbreak{}U/\allowbreak{}Reviews/\allowbreak{}Initial\_\allowbreak{}Independent\_\allowbreak{}Review\_\allowbreak{}Handoff/\allowbreak{}02\_\allowbreak{}INDEPENDENT\_\allowbreak{}REVIEW/\allowbreak{}CHANGES\_\allowbreak{}v0.2\_\allowbreak{}to\_\allowbreak{}v0.3.md}\par
\noindent\textbf{UR}\quad \texttt{C:/\allowbreak{}Users/\allowbreak{}brian/\allowbreak{}Documents/\allowbreak{}Math\ Latex\ etc/\allowbreak{}Papers\ for\ Publication/\allowbreak{}Paper\_\allowbreak{}U\_\allowbreak{}Fresh\_\allowbreak{}Referee\_\allowbreak{}2026-\allowbreak{}09-\allowbreak{}29/\allowbreak{}00\_\allowbreak{}FINAL\_\allowbreak{}REFEREE\_\allowbreak{}REPORT.md}\par
\noindent\textbf{UB}\quad \texttt{C:/\allowbreak{}Users/\allowbreak{}brian/\allowbreak{}Documents/\allowbreak{}Math\ Latex\ etc/\allowbreak{}Papers\ for\ Publication/\allowbreak{}The\ MIND/\allowbreak{}01\_\allowbreak{}Papers/\allowbreak{}02\_\allowbreak{}Hierarchical\_\allowbreak{}Concepts\_\allowbreak{}Paper\_\allowbreak{}U/\allowbreak{}Current\_\allowbreak{}Manuscript/\allowbreak{}v0.3\_\allowbreak{}Rebuilt\_\allowbreak{}Referee\_\allowbreak{}Clarifications/\allowbreak{}Paper\_\allowbreak{}U\_\allowbreak{}Resolution\_\allowbreak{}Limited\_\allowbreak{}Ultrametric\_\allowbreak{}Concepts\_\allowbreak{}v0.3.tex}\par
\noindent\textbf{UF}\quad \texttt{C:/\allowbreak{}Users/\allowbreak{}brian/\allowbreak{}Documents/\allowbreak{}Math\ Latex\ etc/\allowbreak{}Papers\ for\ Publication/\allowbreak{}The\ MIND/\allowbreak{}01\_\allowbreak{}Papers/\allowbreak{}02\_\allowbreak{}Hierarchical\_\allowbreak{}Concepts\_\allowbreak{}Paper\_\allowbreak{}U/\allowbreak{}Supporting\_\allowbreak{}Research/\allowbreak{}Original\_\allowbreak{}Theme\_\allowbreak{}Ultrametrics/\allowbreak{}PAPER\_\allowbreak{}U\_\allowbreak{}RESOLUTION\_\allowbreak{}LIMITED\_\allowbreak{}ULTRAMETRIC\_\allowbreak{}CONCEPTS\_\allowbreak{}DRAFT\_\allowbreak{}v0.1/\allowbreak{}Paper\_\allowbreak{}U\_\allowbreak{}Bridge\_\allowbreak{}to\_\allowbreak{}Paper\_\allowbreak{}F.md}\par
\noindent\textbf{CM}\quad \texttt{C:/\allowbreak{}Users/\allowbreak{}brian/\allowbreak{}Documents/\allowbreak{}Math\ Latex\ etc/\allowbreak{}Papers\ for\ Publication/\allowbreak{}The\ MIND/\allowbreak{}02\_\allowbreak{}Shared\_\allowbreak{}and\_\allowbreak{}Exploratory\_\allowbreak{}Research/\allowbreak{}Measurement\_\allowbreak{}and\_\allowbreak{}RF\_\allowbreak{}IDENT/\allowbreak{}Supporting\_\allowbreak{}Research/\allowbreak{}Height2\_\allowbreak{}Fixed\_\allowbreak{}Library\_\allowbreak{}Gate/\allowbreak{}MIND\_\allowbreak{}TWO\_\allowbreak{}OPERATION\_\allowbreak{}ACCOUNT\_\allowbreak{}TRANSFER\_\allowbreak{}2026-\allowbreak{}09-\allowbreak{}29/\allowbreak{}04\_\allowbreak{}CM\_\allowbreak{}MEASUREMENT\_\allowbreak{}FOUNDATIONS/\allowbreak{}CM\_\allowbreak{}MEASUREMENT\_\allowbreak{}AUDIT.md}\par
\noindent\textbf{T}\quad \texttt{C:/\allowbreak{}Users/\allowbreak{}brian/\allowbreak{}Documents/\allowbreak{}Math\ Latex\ etc/\allowbreak{}Papers\ for\ Publication/\allowbreak{}The\ MIND/\allowbreak{}02\_\allowbreak{}Shared\_\allowbreak{}and\_\allowbreak{}Exploratory\_\allowbreak{}Research/\allowbreak{}Measurement\_\allowbreak{}and\_\allowbreak{}RF\_\allowbreak{}IDENT/\allowbreak{}Original\_\allowbreak{}Theme\_\allowbreak{}Symmetry/\allowbreak{}TWO\_\allowbreak{}OPERATION\_\allowbreak{}CALCULUS\_\allowbreak{}COMPLETE\_\allowbreak{}PLUS\_\allowbreak{}UNIFICATION\_\allowbreak{}2026-\allowbreak{}09-\allowbreak{}29/\allowbreak{}two\_\allowbreak{}operation\_\allowbreak{}calculus/\allowbreak{}TWO\_\allowbreak{}OPERATION\_\allowbreak{}CALCULUS\_\allowbreak{}MASTER\_\allowbreak{}AUDIT.md}\par
\noindent\textbf{TC}\quad \texttt{C:/\allowbreak{}Users/\allowbreak{}brian/\allowbreak{}Documents/\allowbreak{}Math\ Latex\ etc/\allowbreak{}Papers\ for\ Publication/\allowbreak{}The\ MIND/\allowbreak{}02\_\allowbreak{}Shared\_\allowbreak{}and\_\allowbreak{}Exploratory\_\allowbreak{}Research/\allowbreak{}Measurement\_\allowbreak{}and\_\allowbreak{}RF\_\allowbreak{}IDENT/\allowbreak{}Original\_\allowbreak{}Theme\_\allowbreak{}Symmetry/\allowbreak{}TWO\_\allowbreak{}OPERATION\_\allowbreak{}CALCULUS\_\allowbreak{}COMPLETE\_\allowbreak{}PLUS\_\allowbreak{}UNIFICATION\_\allowbreak{}2026-\allowbreak{}09-\allowbreak{}29/\allowbreak{}two\_\allowbreak{}operation\_\allowbreak{}calculus/\allowbreak{}CONSCIOUS\_\allowbreak{}DIFFERENTIATION\_\allowbreak{}STATUS.md}\par
\noindent\textbf{UN}\quad \texttt{C:/\allowbreak{}Users/\allowbreak{}brian/\allowbreak{}Documents/\allowbreak{}Math\ Latex\ etc/\allowbreak{}Papers\ for\ Publication/\allowbreak{}The\ MIND/\allowbreak{}02\_\allowbreak{}Shared\_\allowbreak{}and\_\allowbreak{}Exploratory\_\allowbreak{}Research/\allowbreak{}Measurement\_\allowbreak{}and\_\allowbreak{}RF\_\allowbreak{}IDENT/\allowbreak{}Original\_\allowbreak{}Theme\_\allowbreak{}Symmetry/\allowbreak{}TWO\_\allowbreak{}OPERATION\_\allowbreak{}CALCULUS\_\allowbreak{}COMPLETE\_\allowbreak{}PLUS\_\allowbreak{}UNIFICATION\_\allowbreak{}2026-\allowbreak{}09-\allowbreak{}29/\allowbreak{}two\_\allowbreak{}operation\_\allowbreak{}unification\_\allowbreak{}2026\_\allowbreak{}09\_\allowbreak{}29/\allowbreak{}UNIFIED\_\allowbreak{}DIFFERENTIATION\_\allowbreak{}CALCULUS\_\allowbreak{}INVESTIGATION.md}\par
\noindent\textbf{UA}\quad \texttt{C:/\allowbreak{}Users/\allowbreak{}brian/\allowbreak{}Documents/\allowbreak{}Math\ Latex\ etc/\allowbreak{}Papers\ for\ Publication/\allowbreak{}The\ MIND/\allowbreak{}02\_\allowbreak{}Shared\_\allowbreak{}and\_\allowbreak{}Exploratory\_\allowbreak{}Research/\allowbreak{}Measurement\_\allowbreak{}and\_\allowbreak{}RF\_\allowbreak{}IDENT/\allowbreak{}Supporting\_\allowbreak{}Research/\allowbreak{}Height2\_\allowbreak{}Fixed\_\allowbreak{}Library\_\allowbreak{}Gate/\allowbreak{}MIND\_\allowbreak{}TWO\_\allowbreak{}OPERATION\_\allowbreak{}ACCOUNT\_\allowbreak{}TRANSFER\_\allowbreak{}2026-\allowbreak{}09-\allowbreak{}29/\allowbreak{}07\_\allowbreak{}U1\_\allowbreak{}U6\_\allowbreak{}HOSTILE\_\allowbreak{}AUDIT/\allowbreak{}U1\_\allowbreak{}U6\_\allowbreak{}HOSTILE\_\allowbreak{}AUDIT.md}\par
\noindent\textbf{A1}\quad \texttt{C:/\allowbreak{}Users/\allowbreak{}brian/\allowbreak{}Documents/\allowbreak{}Math\ Latex\ etc/\allowbreak{}Papers\ for\ Publication/\allowbreak{}The\ MIND/\allowbreak{}02\_\allowbreak{}Shared\_\allowbreak{}and\_\allowbreak{}Exploratory\_\allowbreak{}Research/\allowbreak{}Measurement\_\allowbreak{}and\_\allowbreak{}RF\_\allowbreak{}IDENT/\allowbreak{}Supporting\_\allowbreak{}Research/\allowbreak{}Height2\_\allowbreak{}Fixed\_\allowbreak{}Library\_\allowbreak{}Gate/\allowbreak{}MIND\_\allowbreak{}TWO\_\allowbreak{}OPERATION\_\allowbreak{}ACCOUNT\_\allowbreak{}TRANSFER\_\allowbreak{}2026-\allowbreak{}09-\allowbreak{}29/\allowbreak{}07\_\allowbreak{}U1\_\allowbreak{}U6\_\allowbreak{}HOSTILE\_\allowbreak{}AUDIT/\allowbreak{}ADAPTIVE\_\allowbreak{}THEOREM\_\allowbreak{}FINAL.md}\par
\noindent\textbf{H}\quad \texttt{C:/\allowbreak{}Users/\allowbreak{}brian/\allowbreak{}Documents/\allowbreak{}Math\ Latex\ etc/\allowbreak{}Papers\ for\ Publication/\allowbreak{}The\ MIND/\allowbreak{}02\_\allowbreak{}Shared\_\allowbreak{}and\_\allowbreak{}Exploratory\_\allowbreak{}Research/\allowbreak{}Measurement\_\allowbreak{}and\_\allowbreak{}RF\_\allowbreak{}IDENT/\allowbreak{}Supporting\_\allowbreak{}Research/\allowbreak{}Height2\_\allowbreak{}Fixed\_\allowbreak{}Library\_\allowbreak{}Gate/\allowbreak{}research/\allowbreak{}HEIGHT2\_\allowbreak{}ADAPTIVE\_\allowbreak{}MASTER\_\allowbreak{}AUDIT.md}\par
\noindent\textbf{HC}\quad \texttt{C:/\allowbreak{}Users/\allowbreak{}brian/\allowbreak{}Documents/\allowbreak{}Math\ Latex\ etc/\allowbreak{}Papers\ for\ Publication/\allowbreak{}The\ MIND/\allowbreak{}02\_\allowbreak{}Shared\_\allowbreak{}and\_\allowbreak{}Exploratory\_\allowbreak{}Research/\allowbreak{}Measurement\_\allowbreak{}and\_\allowbreak{}RF\_\allowbreak{}IDENT/\allowbreak{}Supporting\_\allowbreak{}Research/\allowbreak{}Height2\_\allowbreak{}Fixed\_\allowbreak{}Library\_\allowbreak{}Gate/\allowbreak{}research/\allowbreak{}COMPLEXITY\_\allowbreak{}AUDIT.md}\par
\noindent\textbf{B}\quad \texttt{C:/\allowbreak{}Users/\allowbreak{}brian/\allowbreak{}Documents/\allowbreak{}Math\ Latex\ etc/\allowbreak{}Papers\ for\ Publication/\allowbreak{}The\ MIND/\allowbreak{}02\_\allowbreak{}Shared\_\allowbreak{}and\_\allowbreak{}Exploratory\_\allowbreak{}Research/\allowbreak{}Measurement\_\allowbreak{}and\_\allowbreak{}RF\_\allowbreak{}IDENT/\allowbreak{}Supporting\_\allowbreak{}Research/\allowbreak{}Bounded\_\allowbreak{}Height2\_\allowbreak{}RF\_\allowbreak{}IDENT/\allowbreak{}THEOREM.md}\par
\noindent\textbf{BR}\quad \texttt{C:/\allowbreak{}Users/\allowbreak{}brian/\allowbreak{}Documents/\allowbreak{}Math\ Latex\ etc/\allowbreak{}Papers\ for\ Publication/\allowbreak{}The\ MIND/\allowbreak{}02\_\allowbreak{}Shared\_\allowbreak{}and\_\allowbreak{}Exploratory\_\allowbreak{}Research/\allowbreak{}Measurement\_\allowbreak{}and\_\allowbreak{}RF\_\allowbreak{}IDENT/\allowbreak{}Reviews/\allowbreak{}Bounded\_\allowbreak{}Height2\_\allowbreak{}Independent\_\allowbreak{}Review/\allowbreak{}REVIEW.md}\par
\noindent\textbf{IDX}\quad \texttt{C:/\allowbreak{}Users/\allowbreak{}brian/\allowbreak{}Documents/\allowbreak{}Math\ Latex\ etc/\allowbreak{}Papers\ for\ Publication/\allowbreak{}The\ MIND/\allowbreak{}00\_\allowbreak{}Index\_\allowbreak{}and\_\allowbreak{}Manifests/\allowbreak{}THEME\_\allowbreak{}CROSSWALK.md}\par
\noindent\textbf{JSON}\quad \texttt{C:/\allowbreak{}Users/\allowbreak{}brian/\allowbreak{}Documents/\allowbreak{}Math\ Latex\ etc/\allowbreak{}Papers\ for\ Publication/\allowbreak{}The\ MIND/\allowbreak{}00\_\allowbreak{}Index\_\allowbreak{}and\_\allowbreak{}Manifests/\allowbreak{}paper\_\allowbreak{}theme\_\allowbreak{}crosswalk.json}\par
\noindent\textbf{MAN}\quad \texttt{C:/\allowbreak{}Users/\allowbreak{}brian/\allowbreak{}Documents/\allowbreak{}Math\ Latex\ etc/\allowbreak{}Papers\ for\ Publication/\allowbreak{}The\ MIND/\allowbreak{}00\_\allowbreak{}Index\_\allowbreak{}and\_\allowbreak{}Manifests/\allowbreak{}source\_\allowbreak{}manifest.csv}\par
\noindent\textbf{ACCESS}\quad \texttt{C:/\allowbreak{}Users/\allowbreak{}brian/\allowbreak{}Documents/\allowbreak{}Math\ Latex\ etc/\allowbreak{}Papers\ for\ Publication/\allowbreak{}The\ MIND/\allowbreak{}00\_\allowbreak{}Index\_\allowbreak{}and\_\allowbreak{}Manifests/\allowbreak{}UNRESOLVED\_\allowbreak{}AND\_\allowbreak{}ACCESS.md}\par
\noindent\textbf{URD}\quad \texttt{C:/\allowbreak{}Users/\allowbreak{}brian/\allowbreak{}Documents/\allowbreak{}Math\ Latex\ etc/\allowbreak{}Papers\ for\ Publication/\allowbreak{}The\ MIND/\allowbreak{}01\_\allowbreak{}Papers/\allowbreak{}02\_\allowbreak{}Hierarchical\_\allowbreak{}Concepts\_\allowbreak{}Paper\_\allowbreak{}U/\allowbreak{}README.md}\par
\endgroup



\end{document}
